% from ega3-1-fr.tex
\oldpage[0\textsubscript{III}]{39}
les limites inductives filtrantes existent et sont exactes, et que $T$ commute
aux limites inductives. Alors $\mathrm R^\bullet T(K^\bullet,K'^\bullet)$ est
un bifoncteur cohomologique en chacun des complexes $K^\bullet$,
$K'^\bullet$~; pour le voir, on se ramène comme dans
(\hyperref[0.11.5.2-fr]{11.5.2}) au cas où $K^\bullet$ et $K'^\bullet$ sont
limités inférieurement~; prenant alors pour $K^\bullet$ et $K'^\bullet$ des
résolutions injectives du type décrit dans
(\hyperref[0.11.5.2.2-fr]{11.5.2.2}), on est ramené à la propriété générale
démontrée dans (M, V, 4.1).
\end{env}

\begin{env}[11.5.6]
\label{0.11.5.6-fr}
De même, les résultats de (\hyperref[0.11.4.7-fr]{11.4.7}) et
(\hyperref[0.11.5.3-fr]{11.5.3}) se généralisent comme suit. Supposons (sous
les hypothèses de (\hyperref[0.11.5.5-fr]{11.5.5})) que l'on ait deux
bicomplexes $L^{\bullet\bullet}=(L^{ij})$,
$L'^{\bullet\bullet}=(L'^{ij})$ tels que $L^{ij}=0$ et $L'^{ij}=0$ pour
$j<0$, que pour tout $i$, $L^{i,\bullet}$ soit une résolution de $K^i$ et
$L'^{i,\bullet}$ une résolution de $K'^i$, et enfin que
$\mathrm R^nT(L^{ij},L'^{hk})=0$ pour $n>0$ et pour tout système d'indices
$(i,j,h,k)$. Alors on a un isomorphisme fonctoriel en $K^\bullet$ et
$K'^\bullet$
\[
  \label{0.11.5.6.1-fr}
  \mathrm R^\bullet T(K^\bullet,K'^\bullet)\xrightarrow{\sim}
  \mathrm H^\bullet(T(L^{\bullet\bullet},L'^{\bullet\bullet})).
  \tag{11.5.6.1}
\]
Cela s'établit comme dans (\hyperref[0.11.5.3-fr]{11.5.3}) en se ramenant au
cas où $K^\bullet$ et $K'^\bullet$ sont limités inférieurement.

Supposons en outre que pour tout couple $(i,j)$ et pour tout couple $(h,k)$,
les foncteurs $A\mapsto T(A,L'^{hk})$ et $A'\mapsto T(L^{ij},A')$ soient
exacts dans $\mathcal C$ et $\mathcal C'$ respectivement. Alors on a aussi des
isomorphismes fonctoriels
\[
  \label{0.11.5.6.2-fr}
  \mathrm R^\bullet T(K^\bullet,K'^\bullet)
  \xrightarrow{\sim}\mathrm H^\bullet(T(L^{\bullet\bullet},K'^\bullet))
  \xrightarrow{\sim}\mathrm H^\bullet(T(K^\bullet,L'^{\bullet\bullet})).
  \tag{11.5.6.2}
\]
La démonstration est semblable à celle de
(\hyperref[0.11.4.7-fr]{11.4.7}).
\end{env}

\begin{env}[11.5.7]
\label{0.11.5.7-fr}
On notera encore que les résultats de (\hyperref[0.11.5.5-fr]{11.5.5}) et
(\hyperref[0.11.5.6-fr]{11.5.6}) sont valables sans supposer que $T$ permute
aux limites inductives, pourvu que l'on se restreigne aux complexes
$K^\bullet$, $K'^\bullet$ limités inférieurement, ou que l'on suppose que tout
objet de $\mathcal C$ (resp. $\mathcal C'$) admet une résolution injective de
longueur bornée, et que dans (\hyperref[0.11.5.6-fr]{11.5.6}) les bicomplexes
$L^{\bullet\bullet}$ et $L'^{\bullet\bullet}$ ont leur second degré limité
supérieurement.
\end{env}

\subsection{Hyperhomologie d'un foncteur par rapport à un complexe
$K_\bullet$}
\label{subsection:0.11.6-fr}

\begin{env}[11.6.1]
\label{0.11.6.1-fr}
Soient $\mathcal C$ une catégorie abélienne,
$K_\bullet=(K_i)_{i\in\mathbb Z}$ un complexe d'objets de $\mathcal C$, dont
l'opérateur de dérivation est de degré $-1$. Une \emph{résolution de
Cartan-Eilenberg gauche} de $K_\bullet$ est formée d'un bicomplexe
$L_{\bullet\bullet}=(L_{ij})$ à opérateurs de dérivation de degré $-1$ avec
$L_{ij}=0$ pour $j<0$, et d'un morphisme de complexes simples
$\varepsilon:L_{\bullet,0}\to K_\bullet$, de façon que les conditions obtenues
à partir de celles de (\hyperref[0.11.4.2-fr]{11.4.2}) par «~renversement des
flèches~» soient remplies. Si tout objet de $\mathcal C$ est quotient d'un
objet projectif, tout complexe $K_\bullet$ de $\mathcal C$ admet une
résolution de Cartan-Eilenberg \emph{projective}, c'est-à-dire formée d'objets
projectifs $L_{ij}$, avec des propriétés fonctorielles semblables à celles de
(\hyperref[0.11.4.2-fr]{11.4.2}). En outre, si $K_\bullet$ est limité
inférieurement (resp. supérieurement), on peut supposer que $L_{ij}=0$ pour
$i<i_0$ (resp. $i>i_0$) si $K_i=0$ pour $i<i_0$ (resp. $i>i_0$). Si tout
objet de $\mathcal C$ admet une résolution projective de longueur $\leq n$,
on peut supposer que $L_{ij}=0$ pour $j>n$.
\end{env}

\begin{env}[11.6.2]
\label{0.11.6.2-fr}
